\chapter{Making Inference Cheaper and Deploying It}\label{ch:efficient}

\Cref{ch:inference} showed that a single decode step of a dense transformer is almost
always \emph{memory-bound}: to emit one token we stream every weight (and the whole KV
cache) from HBM through the compute units, do roughly two FLOPs per byte, and wait.
\Cref{ch:serving} attacked that inefficiency by putting more tokens behind each weight
read: continuous batching, paged KV caches and prefix sharing. This chapter covers the
remaining levers that a production inference stack pulls:
\begin{enumerate}
  \item \textbf{Move fewer bytes.} Store weights (and KV) in 8, 4 or even fewer bits
        (\emph{quantization}) so each decode step reads less memory.
  \item \textbf{Get more tokens per expensive step.} Let a cheap model guess several tokens
        and verify them with \emph{one} forward pass of the big model (\emph{speculative
        decoding}) --- without changing the output distribution at all.
  \item \textbf{Split the model across devices} so that aggregate bandwidth grows with the
        number of GPUs (tensor, pipeline and expert parallel inference).
  \item \textbf{Constrain and ship the output}: structured decoding, OpenAI-compatible
        servers, streaming, load testing --- and, finally, a small browser chat interface
        around our own Qwen3 implementation.
\end{enumerate}
For a systems reader the theme is familiar: every technique trades a little extra compute
(dequantisation arithmetic, draft-model steps, collective communication) for a large
reduction in the scarcest resource, HBM bandwidth per generated token.

\begin{figure}[t]
\centering
\begin{tikzpicture}[
  stage/.style={draw=sysblue,fill=sysbluebg,rounded corners=2pt,minimum width=2.55cm,
    minimum height=1.25cm,align=center,font=\small},
  num/.style={circle,fill=sysblue,text=white,font=\bfseries\small,inner sep=1.5pt},
  arr/.style={-Latex,thick,sysblue!70}]
\node[stage] (s1) {Where a token's\\cost goes};
\node[stage,right=0.45cm of s1] (s2) {Quantization\\(fewer bytes)};
\node[stage,right=0.45cm of s2] (s3) {Speculative\\decoding};
\node[stage,right=0.45cm of s3] (s4) {Parallel \&\\MoE inference};
\node[stage,below=0.9cm of s2] (s5) {Structured\\output};
\node[stage,right=0.45cm of s5] (s6) {Serving stack\\\& load tests};
\node[stage,right=0.45cm of s6,fill=labpurplebg,draw=labpurple] (s7) {Chat interface\\(Chainlit)};
\foreach \i in {1,...,7}{\node[num] at (s\i.north west) {\i};}
\draw[arr] (s1)--(s2); \draw[arr] (s2)--(s3); \draw[arr] (s3)--(s4);
\draw[arr] (s4.south) |- ++(0,-0.45) -| (s5.north);
\draw[arr] (s5)--(s6); \draw[arr] (s6)--(s7);
\node[font=\scriptsize,text=srcgray,below=0.15cm of s1,align=center] {recap of the\\decode roofline};
\node[font=\scriptsize,text=srcgray,below=0.15cm of s7,align=center] {RM App.~G};
\end{tikzpicture}
\caption{Roadmap of this chapter. Stages 2--4 reduce the cost per generated token;
stages 5--7 turn an engine into something people can talk to.}
\label{fig:efficient:roadmap}
\end{figure}

% =====================================================================
\section{Where the cost of a token goes}\label{sec:efficient:cost}

Let a dense decoder have $N$ parameters stored with $b_w$ bytes each, serve a batch of
$B$ sequences whose average context is $T$, and keep $b_{kv}$ bytes per cached element.
With $L$ layers, $G$ KV heads and head dimension $\dhead$, one decode step must read
\begin{equation}
  \text{bytes}_{\text{step}} \;\approx\;
  \underbrace{N\,b_w}_{\text{weights, once per step}} \;+\;
  \underbrace{2\,L\,G\,\dhead\,T\,B\,b_{kv}}_{\text{KV cache of all sequences}}
  \label{eq:efficient:bytes}
\end{equation}
and performs roughly $2NB$ FLOPs for the matrix multiplications plus $4 L H \dhead T B$
for attention scores and value mixing. Dividing by the hardware's bandwidth and peak FLOP
rate gives two lower bounds on the step time; the larger one wins (\cref{ch:pytorch}'s
roofline). With $B=1$ the arithmetic intensity is about one FLOP per byte, two orders of
magnitude below the H100's ridge point of $989/3.35\approx 295$ FLOP/byte.

\begin{napkinmath}[Qwen3 on one H100, batch size 1]
Using \code{llmsys.flops} (\cref{ch:landscape}) and 3.35\,TB/s of HBM bandwidth:
\begin{center}\small
\begin{tabular}{lrrrr}
\toprule
model & params & bf16 weights & ms/token (bf16) & ms/token (int4) \\
\midrule
Qwen3-0.6B & 0.60\,B & 1.19\,GB & 0.36 & 0.09 \\
Qwen3-8B   & 8.19\,B & 16.4\,GB & 4.89 & 1.22 \\
Qwen3-32B  & 32.8\,B & 65.5\,GB & 19.6 & 4.89 \\
\bottomrule
\end{tabular}
\end{center}
These are \emph{floors} from weight traffic alone: 205 tokens/s for Qwen3-8B in bf16,
818 tokens/s if the weights were 4-bit. No kernel can beat them at batch size 1; the
only escapes are fewer bytes (this section), more tokens per byte (batching,
speculation), or more bandwidth (more GPUs).
\end{napkinmath}

The KV term matters as soon as batches and contexts grow. For Qwen3-8B the cache costs
$2\cdot 36\cdot 8\cdot 128\cdot 2$ bytes $=147.5$\,KB per token. A batch of 64 requests
at 4096 tokens of context holds $64\cdot 4096\cdot 147.5\,\text{KB}\approx 38.7$\,GB of
KV, more than twice the weights. \Cref{eq:efficient:bytes} then predicts a 16.4\,ms step
(55\,GB at 3.35\,TB/s), and the step is \emph{still} memory-bound: its arithmetic
intensity is only about 20 FLOP/byte.

\begin{systemslens}[Three levers, three different bottlenecks]
\begin{itemize}
  \item \textbf{Weight quantization} shrinks the first term of \cref{eq:efficient:bytes}.
        It helps most at small batch and short context, where weights dominate.
  \item \textbf{KV quantization} (and GQA, MLA, paging, eviction) shrinks the second term,
        which dominates at large batch or long context, i.e.\ exactly the regime of
        high-throughput serving and long reasoning traces.
  \item \textbf{Speculative decoding} keeps the bytes the same but amortises them over
        several accepted tokens. It helps most when the GPU is idle-ish (small batch)
        and helps little when large batches already make the step compute-bound.
\end{itemize}
Always ask which term of \cref{eq:efficient:bytes} dominates \emph{your} workload before
choosing a technique.
\end{systemslens}

% =====================================================================
\section{Quantization fundamentals}\label{sec:efficient:quant}

\subsection{Number formats}
A quantized model stores tensors in a narrow format and keeps enough metadata (scales,
zero points) to map them back to real numbers. \Cref{fig:efficient:formats} lays out the
formats you will meet. Floating-point formats spend bits on an exponent, which buys
dynamic range: bf16 keeps fp32's 8-bit exponent and is therefore the default training and
inference format. FP8 comes in two flavours: E4M3 (more precision, max $\pm 448$) for
weights and activations, and E5M2 (more range, max $\pm 57344$) traditionally for
gradients. Integer formats spend every bit on resolution but have \emph{uniform} spacing,
so they need an explicit scale to place that grid where the values are.

\begin{figure}[t]
\centering
\begin{tikzpicture}[x=0.36cm,y=0.5cm,font=\scriptsize]
\newcommand{\bits}[5]{% y, sign, exp, mant, label
  \node[anchor=east] at (-0.4,#1+0.5) {#5};
  \fill[warnred!60] (0,#1) rectangle (#2,#1+1);
  \fill[napkin!60] (#2,#1) rectangle (#2+#3,#1+1);
  \fill[sysblue!45] (#2+#3,#1) rectangle (#2+#3+#4,#1+1);
  \foreach \i in {0,...,\numexpr #2+#3+#4-1\relax} \draw[white,line width=0.6pt] (\i,#1) rectangle (\i+1,#1+1);
}
\bits{12}{1}{8}{23}{fp32}
\bits{10.5}{1}{8}{7}{bf16}
\bits{9}{1}{5}{10}{fp16}
\bits{7.5}{1}{4}{3}{fp8 E4M3}
\bits{6}{1}{5}{2}{fp8 E5M2}
\bits{4.5}{1}{2}{1}{fp4 E2M1}
% integers
\node[anchor=east] at (-0.4,3.5) {int8};
\fill[warnred!60] (0,3) rectangle (1,4); \fill[keygreen!45] (1,3) rectangle (8,4);
\foreach \i in {0,...,7} \draw[white,line width=0.6pt] (\i,3) rectangle (\i+1,4);
\node[anchor=east] at (-0.4,2) {int4 / uint4};
\fill[keygreen!45] (0,1.5) rectangle (4,2.5);
\foreach \i in {0,...,3} \draw[white,line width=0.6pt] (\i,1.5) rectangle (\i+1,2.5);
% legend
\fill[warnred!60] (14,3.2) rectangle (14.8,3.8); \node[anchor=west] at (15,3.5) {sign};
\fill[napkin!60] (18,3.2) rectangle (18.8,3.8); \node[anchor=west] at (19,3.5) {exponent};
\fill[sysblue!45] (23.5,3.2) rectangle (24.3,3.8); \node[anchor=west] at (24.5,3.5) {mantissa};
\fill[keygreen!45] (14,1.7) rectangle (14.8,2.3); \node[anchor=west] at (15,2) {integer magnitude (uniform grid, needs a scale)};
\end{tikzpicture}
\caption{Bit layouts of the formats used for LLM inference. Exponent bits buy dynamic
range; mantissa or integer bits buy resolution. Formats narrower than 8 bits are only
usable together with fine-grained (per-group or per-block) scales.}
\label{fig:efficient:formats}
\end{figure}

\subsection{Uniform affine quantization}
A $b$-bit uniform quantizer maps a real tensor $x$ to integer codes with a scale $s>0$
and a zero point $z$:
\begin{equation}
  q = \operatorname{clamp}\!\Big(\big\lfloor x/s \big\rceil + z,\; q_{\min},\; q_{\max}\Big),
  \qquad \hat x = s\,(q - z).
  \label{eq:efficient:affine}
\end{equation}
\textbf{Symmetric} (absmax) quantization sets $z=0$ and $s=\max|x|/(2^{b-1}-1)$, which
suits roughly zero-centred weights and makes the integer GEMM cheap (no zero-point
correction terms). \textbf{Asymmetric} quantization uses the full range
$[\min x,\max x]$ with $s=(\max x-\min x)/(2^b-1)$ and $z=\lfloor -\min x/s\rceil$, which
wastes no codes on one-sided distributions.

If values are spread across the grid, the rounding error is roughly uniform on
$[-s/2,s/2]$ and has variance $s^2/12$. Each extra bit halves $s$, which gives the
familiar rule of thumb of about 6\,dB of signal-to-quantization-noise ratio per bit.
The catch is the word \emph{spread}: one outlier stretches $\max|x|$, inflates $s$ for
\emph{everything else} in the same group, and pushes small values to zero. Nearly every
modern quantization method is a way of fighting outliers.

\begin{lstlisting}[caption={Symmetric absmax and asymmetric group-wise quantization
(from \ident{llmsys/quant.py}).},label={lst:efficient:absmax}]
def quantize_absmax(w, bits=8, dim=-1):
    qmax = 2 ** (bits - 1) - 1
    amax = w.abs().max() if dim is None else w.abs().amax(dim=dim, keepdim=True)
    scale = (amax / qmax).clamp(min=1e-12).float()
    q = torch.clamp(torch.round(w.float() / scale), -qmax, qmax).to(torch.int8)
    return q, scale

def quantize_groupwise(w, bits=4, group_size=128):
    out_f, in_f = w.shape
    qmax = 2 ** bits - 1
    g = w.float().reshape(out_f, in_f // group_size, group_size)
    wmin, wmax = g.amin(dim=-1), g.amax(dim=-1)
    scale = ((wmax - wmin) / qmax).clamp(min=1e-12)
    zero = torch.round(-wmin / scale).clamp(0, qmax)
    q = torch.clamp(torch.round(g / scale[..., None]) + zero[..., None], 0, qmax)
    return q.to(torch.uint8).reshape(out_f, in_f), scale, zero
\end{lstlisting}

\subsection{Granularity: per-tensor, per-channel, per-group}
\Cref{fig:efficient:granularity} shows the three granularities for a weight matrix
$W\in\R^{d_{\text{out}}\times d_{\text{in}}}$ (one row per output channel, as stored by
\code{nn.Linear}). Per-tensor uses one scale and is fragile. Per-output-channel uses one
scale per row and is the default for int8 weights. Since $y_j=\sum_i W_{ji}x_i$, a
per-row scale factors out of the dot product and can be applied once to the GEMM
result. Per-group splits each row into groups of $g$ consecutive input channels
(typically $g=128$, sometimes 32 or 64), each with its own scale and zero point. This is
what makes 4-bit weights usable.

\begin{figure}[t]
\centering
\resizebox{\linewidth}{!}{%
\begin{tikzpicture}[x=0.28cm,y=0.28cm]
\pgfmathsetseed{7}
\newcommand{\wgrid}[2]{% xoffset, title
  \begin{scope}[xshift=#1 cm]
  \foreach \r in {0,...,7}{\foreach \c in {0,...,15}{
     \pgfmathsetmacro{\v}{10+40*rnd}
     \fill[sysblue!\v] (\c,-\r) rectangle ++(1,-1);}}
  \draw[gray] (0,0) rectangle (16,-8);
  \node[font=\small\bfseries,above] at (8,0.2) {#2};
  \end{scope}}
\wgrid{0}{per-tensor}
\draw[warnred,very thick] (0,0) rectangle (16,-8);
\node[font=\scriptsize,below,align=center] at (8,-8.3) {1 scale\\outlier hurts all};
\wgrid{5.6}{per-channel}
\foreach \r in {0,...,7}{\draw[warnred,thick] (20,-\r) rectangle (36,-\r-1);}
\node[font=\scriptsize,below,align=center] at (28,-8.3) {$d_{\text{out}}$ scales\\(one per row)};
\wgrid{11.2}{per-group ($g=4$)}
\foreach \r in {0,...,7}{\foreach \c in {0,4,8,12}{\draw[warnred,thick] (40+\c,-\r) rectangle (44+\c,-\r-1);}}
\node[font=\scriptsize,below,align=center] at (48,-8.3) {$d_{\text{out}}\cdot d_{\text{in}}/g$ scales\\(+ zero points)};
\draw[-Latex,thick] (-1.2,-0.5) -- (-1.2,-7.5) node[midway,left,font=\scriptsize] {out};
\draw[-Latex,thick] (0.5,-11.6) -- (15.5,-11.6) node[midway,below,font=\scriptsize] {in};
\end{tikzpicture}}
\caption{Quantization granularity for a weight matrix. Red boxes enclose the elements
that share one scale. Finer groups track local magnitudes better at the cost of more
metadata.}
\label{fig:efficient:granularity}
\end{figure}

\begin{napkinmath}[The metadata tax of group-wise int4]
With $g=128$, an fp16 scale and a 4-bit zero point per group cost $(16+4)/128\approx0.16$
extra bits per weight; with an fp16 zero point, $32/128=0.25$ bits. So ``4-bit'' Qwen3-8B
actually stores about $4.16$--$4.25$ bits per parameter: roughly 4.3\,GB instead of the
naive 4.1\,GB. With $g=32$ the tax rises to 0.5--1 bit, which is why tiny groups are
reserved for formats with very cheap (8-bit exponent-only) scales such as the
microscaling (MX) formats.
\end{napkinmath}

\subsection{Packing sub-byte values}
Memory is byte-addressed, so two 4-bit codes share one byte. Our \code{pack\_int4}
places even columns in the low nibble and odd columns in the high nibble
(\cref{fig:efficient:pack}). Production kernels choose an interleaving that matches how
their GEMM loads registers. The Marlin kernel, for instance, pre-shuffles weights
offline so that one 128-bit load feeds tensor-core fragments directly. Packing order is
therefore part of the \emph{kernel's} contract, not the model's, and checkpoints are
often repacked at load time.

\begin{figure}[t]
\centering
\begin{tikzpicture}[font=\small,
  c/.style={draw,minimum width=0.85cm,minimum height=0.6cm,fill=keygreen!25},
  b/.style={draw,minimum width=1.7cm,minimum height=0.6cm,fill=sysblue!20}]
\foreach \i/\v in {0/3,1/12,2/7,3/0,4/15,5/9}{\node[c] (q\i) at (\i*0.85,0) {\v};
  \node[font=\scriptsize,text=srcgray] at (\i*0.85,0.55) {$q_{\i}$};}
\foreach \j/\lo/\hi in {0/3/12,1/7/0,2/15/9}{
  \node[b] (p\j) at (\j*1.7+0.425,-1.9) {};}
\node[b] at (0.425,-1.9) {\texttt{0xC3}};
\node[b] at (2.125,-1.9) {\texttt{0x07}};
\node[b] at (3.825,-1.9) {\texttt{0x9F}};
\foreach \j in {0,1,2}{
  \pgfmathtruncatemacro{\a}{2*\j}\pgfmathtruncatemacro{\b}{2*\j+1}
  \draw[-Latex] (q\a.south) -- ($(p\j.north)+(-0.4,0)$);
  \draw[-Latex] (q\b.south) -- ($(p\j.north)+(0.4,0)$);}
\node[anchor=west,font=\scriptsize,align=left] at (5.3,-1.9)
  {byte $= q_{2i}\;|\;(q_{2i+1}\ll 4)$\\unpack: \texttt{lo = p \& 0xF}, \texttt{hi = p >> 4}};
\node[anchor=west,font=\scriptsize] at (5.3,0) {6 codes $\times$ 4 bits};
\end{tikzpicture}
\caption{Packing pairs of 4-bit codes into bytes: the low nibble holds the even
element and the high nibble the odd element.}
\label{fig:efficient:pack}
\end{figure}

\subsection{Weight-only quantization (W8A16, W4A16)}\covers{RM D.5}
In weight-only schemes the activations stay in bf16 and the weights are dequantized
\emph{inside} the GEMM: a kernel loads packed integers from HBM, expands them in
registers, multiplies by the group scale and feeds tensor cores in bf16. Because decode
is memory-bound, the extra ALU work is free and the step time drops nearly in proportion
to the weight bytes. \Cref{lst:efficient:quantlinear} is the reference layer we use in
the lab. It dequantizes the full weight in PyTorch, which is correct but \emph{not}
faster: it materialises the fp32 weight in memory, the very traffic we wanted to avoid.
\Cref{app:triton} writes the fused kernel.

\begin{lstlisting}[caption={A weight-only quantized linear layer (reference semantics,
not a fast kernel).},label={lst:efficient:quantlinear}]
class QuantLinear(nn.Module):
    def __init__(self, linear, bits=8, group_size=128):
        super().__init__()
        w = linear.weight.data
        if bits == 8:                       # per-output-channel symmetric
            q, s = quantize_absmax(w, 8, dim=-1)
            self.register_buffer("qweight", q); self.register_buffer("scale", s)
        else:                               # 4-bit asymmetric groups, packed
            q, s, z = quantize_groupwise(w, 4, group_size)
            self.register_buffer("qweight", pack_int4(q))
            self.register_buffer("scale", s.half()); self.register_buffer("zero", z)
        ...
    def forward(self, x):
        w = self.dequantized_weight(x.dtype)   # a fused kernel never builds this
        return x @ w.T + (0 if self.bias is None else self.bias)
\end{lstlisting}

\begin{pitfall}[Quantizing the wrong layers]
Embeddings and the output head are often kept in bf16. The head maps to a vocabulary of
about 150k entries, so its errors land directly on the next-token distribution. Our
\code{quantize\_model} skips \code{out\_head} by default. Norm weights and biases are a
negligible fraction of bytes; never quantize them.
\end{pitfall}

\subsection{GPTQ: rounding with error feedback}
Round-to-nearest (RTN) treats each weight independently. But what we care about is the
layer's \emph{output} on realistic inputs. Given calibration activations
$X\in\R^{n\times d_{\text{in}}}$ for a layer, GPTQ (Frantar et al., 2023) minimises the
output error row by row:
\begin{equation}
  \min_{Q}\; \big\| X W^\top - X Q^\top \big\|_F^2
  \;=\; \sum_j (\vect w_j - \vect q_j)^\top H\, (\vect w_j - \vect q_j),
  \qquad H = X^\top X .
  \label{eq:efficient:gptq}
\end{equation}
It builds on Optimal Brain Surgeon: quantize one input column $i$, and update the
remaining unquantized columns of every row to absorb the error, using the inverse
Hessian:
\begin{equation}
  \vect w_{:,>i} \;\leftarrow\; \vect w_{:,>i} \;-\;
  \frac{w_{:,i}-q_{:,i}}{[H^{-1}]_{ii}}\; [H^{-1}]_{i,>i} .
\end{equation}
GPTQ's engineering contributions make this fast enough for 100B-parameter models. It
quantizes columns in a fixed order, which makes $H^{-1}$ the same for every row. It
uses a Cholesky factorisation of $H^{-1}$, whose rows provide the update coefficients.
And it uses lazy block updates, which turn many rank-1 updates into GEMMs.
\Cref{lst:efficient:gptq} is the unblocked version.

\begin{lstlisting}[caption={GPTQ without lazy batching. On correlated calibration
inputs it gives lower output error than RTN on the same grid.},label={lst:efficient:gptq}]
def gptq_quantize(w, x, bits=4, damp=0.01):
    W = w.clone().float(); d = W.shape[1]
    H = x.float().T @ x.float()
    H += damp * torch.mean(torch.diag(H)) * torch.eye(d)        # dampening
    Hinv = torch.linalg.cholesky(
        torch.cholesky_inverse(torch.linalg.cholesky(H)), upper=True)
    qmax = 2 ** (bits - 1) - 1
    scale = (W.abs().amax(dim=1) / qmax).clamp(min=1e-12)      # per-row grid
    Q = torch.zeros_like(W)
    for i in range(d):
        col = W[:, i]
        q = torch.clamp(torch.round(col / scale), -qmax, qmax) * scale
        Q[:, i] = q
        err = (col - q) / Hinv[i, i]
        W[:, i + 1:] -= err[:, None] * Hinv[i, i + 1:][None, :]  # error feedback
    return Q
\end{lstlisting}

AWQ (Lin et al., 2024) starts from a different observation: a small fraction of weight
input channels matter far more because their \emph{activations} are large. Instead of
keeping those channels in high precision (which breaks the uniform kernel), AWQ
multiplies salient input channels of $W$ by a factor $s>1$ and divides the
corresponding activations by $s$. The product is unchanged, but the important weights
now occupy more of the quantization grid. The factors are searched on calibration data.
Both GPTQ and AWQ produce ordinary group-wise int4 checkpoints that the same W4A16
kernels can serve.

\subsection{Quantizing activations: outliers and SmoothQuant}
Weight-only schemes do nothing for compute-bound prefill or large-batch decode. There
we want int8 (or FP8) \emph{matmuls}, which means quantizing activations too (W8A8).
LLM activations are hard. Beyond a few billion parameters, a handful of hidden
dimensions carry values 10--100$\times$ larger than the rest, consistently across
tokens. LLM.int8() (Dettmers et al., 2022) handles them with a mixed-precision
decomposition: outlier columns in fp16, the rest in int8. SmoothQuant (Xiao et al.,
2023) instead migrates the difficulty into the weights with a per-input-channel factor:
\begin{equation}
  Y = X W^\top = \big(X\,\mathrm{diag}(\vect s)^{-1}\big)\big(W\,\mathrm{diag}(\vect s)\big)^\top,
  \qquad s_j = \frac{\max|X_{:,j}|^{\alpha}}{\max|W_{:,j}|^{1-\alpha}} .
  \label{eq:efficient:smooth}
\end{equation}
With $\alpha\approx0.5$ the per-channel ranges of activations and weights are balanced.
The division by $\vect s$ is folded offline into the preceding RMSNorm's scale, so it
costs nothing at runtime (\cref{fig:efficient:smooth}).

\begin{figure}[t]
\centering
\resizebox{\linewidth}{!}{%
\begin{tikzpicture}[x=0.42cm,y=0.06cm,font=\scriptsize]
\newcommand{\bars}[3]{% xshift, list, color
 \begin{scope}[xshift=#1 cm]
 \foreach \h [count=\i] in {#2} {\fill[#3] (\i-1,0) rectangle (\i-0.25,\h);}
 \draw[gray] (-0.3,0) -- (8,0);
 \end{scope}}
\bars{0}{4,5,60,3,6,4,55,5}{warnred!60}
\node[font=\small] at (3.8,70) {$\max|X_{:,j}|$};
\bars{4.2}{6,5,7,6,4,6,5,7}{sysblue!50}
\node[font=\small] at (13.8,70) {$\max|W_{:,j}|$};
\draw[-Latex,very thick] (19.5,30) -- (22.5,30) node[midway,above] {$\div\vect s$, $\times\vect s$};
\bars{10.0}{11,12,30,10,12,11,28,12}{warnred!60}
\node[font=\small] at (27.8,70) {$\max|X'_{:,j}|$};
\bars{14.2}{12,11,24,11,10,12,22,13}{sysblue!50}
\node[font=\small] at (37.8,70) {$\max|W'_{:,j}|$};
\node[align=center] at (31.8,-12) {both now quantize well with one scale per tensor/channel};
\node[align=center] at (7.8,-12) {two outlier channels dominate the activation scale};
\end{tikzpicture}}
\caption{SmoothQuant moves per-channel magnitude from the activations into the weights.
The layer computes the same function, but both operands become int8-friendly.}
\label{fig:efficient:smooth}
\end{figure}

\subsection{FP8 and microscaling formats}
Hopper-class GPUs execute FP8 GEMMs at twice the bf16 rate, and Blackwell adds FP4.
FP8 E4M3 has enough dynamic range that a single per-tensor scale (computed from the
running absmax, possibly from the previous iteration: ``delayed scaling'') is often
sufficient for W8A8 inference with negligible quality loss. Below 8 bits even floating
point needs block scales. The OCP microscaling formats (MXFP8/MXFP6/MXFP4) attach one
shared power-of-two (E8M0) scale to every block of 32 elements. NVIDIA's NVFP4 uses
blocks of 16 with an FP8 scale plus a per-tensor scale. Hardware applies the block
scales inside the tensor-core instruction, so the fine granularity no longer costs a
separate dequantization pass.

\begin{lstlisting}[caption={Simulating FP8 storage with a per-tensor scale.},
label={lst:efficient:fp8}]
def fp8_roundtrip(t, fmt="e4m3"):
    dtype = {"e4m3": torch.float8_e4m3fn, "e5m2": torch.float8_e5m2}[fmt]
    scale = (t.abs().max() / torch.finfo(dtype).max).clamp(min=1e-12)
    return (t / scale).to(dtype).to(torch.float32) * scale
\end{lstlisting}

\subsection{Quantizing the KV cache}
\Cref{eq:efficient:bytes} shows that the KV cache dominates long-context and
large-batch decode, so quantizing it pays twice: fewer bytes per step, and more
sequences fit in memory, which raises the achievable batch size. Values behave like
ordinary activations and quantize well per token. Keys are different. Empirically, a few \emph{channels} of $K$ are consistently large
across tokens, so KIVI (Liu et
al., 2024) quantizes keys per channel and values per token. FP8 KV caches are a common
production default; int4 and 2-bit caches need care, usually by keeping a recent window
in full precision. Since each attention kernel reads the cache, the dequantization must
be fused into the attention kernel as well (\cref{app:triton}).

\begin{napkinmath}[Which bytes to cut for a throughput server]
Qwen3-8B, batch 64, 4096-token contexts, one H100 (numbers from
\cref{sec:efficient:cost}):
\begin{center}\small
\begin{tabular}{lrrr}
\toprule
configuration & weights & KV cache & step floor \\
\midrule
bf16 weights, bf16 KV & 16.4\,GB & 38.7\,GB & 16.4\,ms \\
int4 weights, bf16 KV & 4.1\,GB  & 38.7\,GB & 12.8\,ms \\
bf16 weights, fp8 KV  & 16.4\,GB & 19.3\,GB & 10.7\,ms \\
int4 weights, fp8 KV  & 4.1\,GB  & 19.3\,GB & 7.0\,ms \\
\bottomrule
\end{tabular}
\end{center}
In this regime KV quantization alone beats 4-bit weights. And the first row would not
even fit next to activations and CUDA graphs on an 80\,GB card.
\end{napkinmath}

\subsection{Did quantization break the model?}\covers{RM D.5}
Measure the quantized model against the reference model, not just against a benchmark
score:
\begin{itemize}
  \item \textbf{Per-token KL divergence} $\KL(p_{\text{ref}}\,\|\,p_{\text{quant}})$ on held-out
        text. It is sensitive, cheap, and needs no labels.
  \item \textbf{Perplexity delta}, the classic metric. It hides catastrophic errors on
        rare tokens.
  \item \textbf{Task accuracy with generation} (MATH-500, \cref{ch:eval}). Reasoning
        models are especially sensitive: small per-token errors compound over traces of
        thousands of tokens, and a flipped digit is a wrong answer.
  \item \textbf{Flip rate}: the fraction of benchmark items whose correctness changed.
        Equal average accuracy can hide many flips in both directions.
\end{itemize}

\begin{deepdive}[Larger models: quantize or go smaller?]
For a fixed memory budget, a larger model at 4 bits usually beats a smaller model at 16
bits. This was a central finding of the ``k-bit inference scaling laws'' study (Dettmers
\& Zettlemoyer, 2023), whose optimum sat around 4 bits. Practically: if Qwen3-8B in bf16
does not fit your card but Qwen3-14B in int4 (about 7.5\,GB of weights) does, try the
latter first. Below about 3 bits, quality typically falls off a cliff unless the method
is specialised (vector quantization, quantization-aware fine-tuning).
\end{deepdive}

% =====================================================================
\section{Speculative decoding}\label{sec:efficient:spec}

At batch size 1, verifying five tokens costs about as much as generating one. The
forward pass over a $(1\times 5)$ block reads the same weights as a $(1\times1)$ step,
and the extra FLOPs are invisible below the roofline ridge. Speculative decoding (Leviathan
et al., 2023; Chen et al., 2023) exploits this. A small \emph{draft} model $q$
proposes $\gamma$ tokens autoregressively (cheap steps). The large \emph{target} model
$p$ then scores all of them in \emph{one} forward pass, and a rejection-sampling rule
keeps the longest acceptable prefix plus one extra token
(\cref{fig:efficient:spec-timeline}).

\begin{figure}[t]
\centering
\resizebox{\linewidth}{!}{%
\begin{tikzpicture}[font=\scriptsize,
  dstep/.style={draw=keygreen,fill=keygreen!20,minimum width=0.55cm,minimum height=0.5cm},
  tstep/.style={draw=sysblue,fill=sysblue!25,minimum height=0.5cm},
  tok/.style={minimum width=0.62cm,minimum height=0.5cm,draw,font=\scriptsize\ttfamily}]
% baseline
\node[anchor=east,font=\small] at (-0.2,2.2) {target only};
\foreach \i in {0,...,6}{\node[tstep,minimum width=1.5cm] at (0.8+\i*1.6,2.2) {$p$ step};}
\node[anchor=west] at (11.7,2.2) {7 tokens, 7 big steps};
% speculative
\node[anchor=east,font=\small] at (-0.2,0.6) {speculative};
\foreach \i in {0,...,3}{\node[dstep] at (0.3+\i*0.6,0.6) {$q$};}
\node[tstep,minimum width=1.5cm] (v1) at (3.45,0.6) {verify 5};
\foreach \i in {0,...,3}{\node[dstep] at (4.5+\i*0.6,0.6) {$q$};}
\node[tstep,minimum width=1.5cm] (v2) at (7.65,0.6) {verify 5};
\node[anchor=west] at (8.6,0.6) {7 tokens, 2 big + 8 small steps};
% round detail
\node[tok,fill=keygreen!20] (d1) at (0.6,-0.8) {the};
\node[tok,fill=keygreen!20,right=0.05cm of d1] (d2) {cat};
\node[tok,fill=keygreen!20,right=0.05cm of d2] (d3) {sat};
\node[tok,fill=warnred!25,right=0.05cm of d3] (d4) {in};
\node[tok,fill=keygreen!8,right=0.05cm of d4,text=gray] (d5) {the};
\node[anchor=west,right=0.15cm of d5] (arrow) {$\Rightarrow$};
\node[tok,fill=keygreen!20,right=0.15cm of arrow] (a1) {the};
\node[tok,fill=keygreen!20,right=0.05cm of a1] (a2) {cat};
\node[tok,fill=keygreen!20,right=0.05cm of a2] (a3) {sat};
\node[tok,fill=sysblue!25,right=0.05cm of a3] (a4) {on};
\node[anchor=west,right=0.2cm of a4,align=left] {accepted 3 drafts; rejected ``in'',\\resampled ``on'' from $\mathrm{norm}(\max(p-q,0))$};
\node[above=0.05cm of d2,font=\scriptsize\itshape,text=srcgray] {round 1 drafts ($\gamma=4$)};
\end{tikzpicture}}
\caption{Speculative decoding timeline. In each round the draft model takes $\gamma$
cheap steps, then the target verifies the pending token plus all $\gamma$ drafts in a
single forward pass. Every round commits between 1 and $\gamma+1$ tokens.}
\label{fig:efficient:spec-timeline}
\end{figure}

\subsection{The acceptance rule and why it is exact}
Let $p(\cdot)$ and $q(\cdot)$ be the target and draft distributions at some position
(after temperature and top-$k$/top-$p$ filtering, the distributions we would actually
sample from). The draft proposes $x\sim q$. We accept it with probability
$\min\!\big(1, p(x)/q(x)\big)$. If we reject, we sample a replacement from the
\emph{residual} distribution
\begin{equation}
  r(x) = \frac{\max\big(0,\,p(x)-q(x)\big)}{\sum_y \max\big(0,\,p(y)-q(y)\big)} .
\end{equation}

\begin{proposition}[Speculative sampling is lossless]\label{prop:efficient:exact}
The token emitted by this procedure is distributed exactly according to $p$.
\end{proposition}
\begin{proof}
The probability of accepting is
$\beta=\sum_y q(y)\min(1,p(y)/q(y))=\sum_y \min(p(y),q(y))$. Since
$\sum_y p(y)=1$, the residual's normaliser is
\[
\textstyle\sum_y \max(0,p(y)-q(y)) = \sum_y \big(p(y)-\min(p(y),q(y))\big) = 1-\beta .
\]
Hence
\begin{align*}
\Pr[\text{emit } x] &= \underbrace{q(x)\min\!\Big(1,\tfrac{p(x)}{q(x)}\Big)}_{\text{draft accepted}}
\;+\; \underbrace{(1-\beta)\,\frac{\max(0,p(x)-q(x))}{1-\beta}}_{\text{rejected, then resampled}}\\
&= \min(p(x),q(x)) + \max(0,p(x)-q(x)) \;=\; p(x). \qedhere
\end{align*}
\end{proof}

Applied position by position, the proposition shows that the whole generated sequence
has the target's distribution. Each later draft was proposed conditioned on accepted
predecessors, and verification stops at the first rejection. With temperature 0 both
distributions are one-hot. The rule then reduces to ``accept while the draft equals the
target's argmax, then emit the target's argmax'', so greedy speculative decoding
reproduces greedy target decoding \emph{token for token}. The lab tests exactly this.
Note that $\beta = 1-\mathrm{TV}(p,q)$: the acceptance rate is one minus the total
variation distance between the two models.

\begin{lstlisting}[caption={The accept-or-resample rule (\ident{llmsys/specdec.py}).},
label={lst:efficient:accept}]
def accept_or_resample(p, q, draft_token, generator=None):
    px, qx = float(p[draft_token]), float(q[draft_token])
    u = torch.rand((), generator=generator).item()
    if qx > 0 and u < min(1.0, px / qx):
        return True, draft_token
    residual = torch.clamp(p - q, min=0)
    s = residual.sum()
    residual = p if s <= 0 else residual / s
    return False, int(torch.multinomial(residual, 1, generator=generator))
\end{lstlisting}

\subsection{The full loop with KV caches and rollback}
The systems content of speculative decoding is mostly cache management. We keep both
models' KV caches covering ``all committed tokens except the last''. The last committed
token is \emph{pending}: it is the next input to feed. A round then proceeds as follows:
\begin{enumerate}
  \item Feed the pending token to the draft, sample $d_1$; feed $d_1$, sample $d_2$; \dots;
        sample $d_\gamma$. The draft cache grows by $\gamma$ entries ($d_\gamma$ itself is not fed).
  \item Feed the block $[\text{pending}, d_1,\dots,d_\gamma]$ to the target in one call. Its
        $\gamma+1$ output rows are the distributions $p_1,\dots,p_{\gamma+1}$.
  \item Accept a prefix $d_1..d_k$ and append one extra token: the residual sample at the
        first rejection, or a ``bonus'' sample from $p_{\gamma+1}$ if all were accepted.
  \item Roll both caches back by truncation to the new committed length minus one. If all
        drafts were accepted, the draft model never saw $d_\gamma$, so feed it before the
        next round.
\end{enumerate}

\begin{lstlisting}[caption={One round of speculative decoding (abridged from
\ident{speculative_generate}).},label={lst:efficient:specloop}]
nxt = torch.tensor([[seq[-1]]])                      # pending token
for _ in range(g):                                   # 1) draft proposes
    q = probs_from_logits(draft(nxt, cache=dc)[0, -1], temperature, top_k, top_p)
    tok = int(torch.multinomial(q, 1, generator=generator))
    drafts.append(tok); qs.append(q); nxt = torch.tensor([[tok]])
block = torch.tensor([[seq[-1]] + drafts])           # 2) one target pass
ps = probs_from_logits(target(block, cache=tc)[0], temperature, top_k, top_p)
k, extra = 0, None                                   # 3) verify
for i, tok in enumerate(drafts):
    ok, t = accept_or_resample(ps[i], qs[i], tok, generator)
    if not ok:
        extra = t; break
    k += 1
if extra is None:
    extra = int(torch.multinomial(ps[g], 1, generator=generator))
seq.extend(drafts[:k] + [extra])                     # 4) commit + roll back
truncate_cache(tc, len(seq) - 1)
\end{lstlisting}

\begin{pitfall}[Rollback in paged and static caches]
Truncating our concatenation-based \code{KVCache} is a slice. In a paged engine
(\cref{ch:serving}) rollback means decrementing the sequence length and possibly freeing
the last block. In a static cache it only means resetting a length counter. Rejected
entries are simply overwritten next round. Engines that capture the decode step in a CUDA
graph must also capture a separate ``verify $\gamma+1$ tokens'' graph, because shapes
differ.
\end{pitfall}

\subsection{How many tokens per round, and how much faster?}
Suppose each draft is accepted independently with probability $\alpha$. The number of
committed tokens per round is one plus a truncated geometric variable:
\begin{equation}
  \E[\text{tokens per round}] \;=\; \sum_{i=0}^{\gamma}\alpha^i \;=\;
  \frac{1-\alpha^{\gamma+1}}{1-\alpha}.
  \label{eq:efficient:tpr}
\end{equation}
Let $c$ be the draft-step time relative to a target step, and assume verifying $\gamma+1$
tokens costs one target step (memory-bound). The wall-clock speedup is then
\begin{equation}
  S(\alpha,\gamma,c) \;=\; \frac{1-\alpha^{\gamma+1}}{(1-\alpha)(\gamma c + 1)} .
  \label{eq:efficient:speedup}
\end{equation}
\Cref{fig:efficient:spec-curves} plots both. Tokens per round saturate at
$1/(1-\alpha)$, while the cost of drafting grows linearly in $\gamma$, so every
$(\alpha,c)$ pair has an optimal $\gamma$, typically 3--8 for a good draft.

\begin{figure}[t]
\centering
\begin{tikzpicture}
\begin{groupplot}[group style={group size=2 by 1,horizontal sep=1.6cm},
  width=0.47\linewidth,height=5cm,xlabel={$\gamma$ (draft length)},
  xmin=0,xmax=12,grid=major,grid style={gray!20},legend style={font=\scriptsize},
  every axis plot/.append style={thick,mark=none,domain=0:12,samples=13}]
\nextgroupplot[ylabel={tokens per target call},legend pos=north west,ymin=0.8,ymax=10.5]
\addplot[sysblue] {(1-0.6^(x+1))/(1-0.6)}; \addlegendentry{$\alpha=0.6$}
\addplot[keygreen] {(1-0.8^(x+1))/(1-0.8)}; \addlegendentry{$\alpha=0.8$}
\addplot[warnred] {(1-0.9^(x+1))/(1-0.9)}; \addlegendentry{$\alpha=0.9$}
\nextgroupplot[ylabel={speedup $S$ (draft cost $c=0.05$)},legend pos=north east,ymin=0.8,ymax=5]
\addplot[sysblue] {(1-0.6^(x+1))/((1-0.6)*(0.05*x+1))}; \addlegendentry{$\alpha=0.6$}
\addplot[keygreen] {(1-0.8^(x+1))/((1-0.8)*(0.05*x+1))}; \addlegendentry{$\alpha=0.8$}
\addplot[warnred] {(1-0.9^(x+1))/((1-0.9)*(0.05*x+1))}; \addlegendentry{$\alpha=0.9$}
\end{groupplot}
\end{tikzpicture}
\caption{Analytic model of speculative decoding (\cref{eq:efficient:tpr,eq:efficient:speedup}).
Left: committed tokens per target forward pass. Right: speedup when a draft step costs 5\%
of a target step. Each curve has an interior optimum.}
\label{fig:efficient:spec-curves}
\end{figure}

\begin{napkinmath}[Qwen3-0.6B drafting for Qwen3-8B]
Both models share a tokenizer, a hard requirement for vanilla speculative decoding. At
batch 1 their weight-read floors are 0.36\,ms and 4.89\,ms, so $c\approx 0.07$. With an
acceptance rate of $\alpha=0.8$ and $\gamma=4$:
$\E[\text{tokens}]=(1-0.8^5)/0.2=3.36$ and the round costs $4(0.07)+1=1.28$ target
steps, giving $S\approx 2.6\times$. Real systems see less: the 0.6B model's step is far
from its bandwidth floor (kernel-launch overheads dominate tiny models, \cref{ch:inference}),
which makes the effective $c$ larger, and $\alpha$ varies wildly by domain. Code and math
are often highly predictable; creative writing at high temperature is not.
\end{napkinmath}

\subsection{Measuring acceptance correctly}
It is tempting to report ``accepted drafts / proposed drafts''. That is a \emph{biased}
estimate of $\alpha$: after the first rejection in a round, the remaining drafts are
never examined, yet they are still counted as proposed. A consistent estimator counts
only examined positions,
$\hat\alpha=\text{accepted}/(\text{accepted}+\text{rejections})$, and this is what
\code{SpecStats.alpha\_hat} reports.

We measured both on CPU with toy models in this environment. The target is a 4-layer
random Qwen3-shaped model (untied embeddings), sampled at temperature 0.25. The drafts
are copies of the target with Gaussian noise added to every weight, so more noise means
a worse draft. The run generates 160 tokens per configuration
(\code{bench.py toy} in the lab):
\begin{center}\small
\begin{tabular}{rrrrrr}
\toprule
draft noise & $\gamma$ & kept/proposed & $\hat\alpha$ & tokens/target call & \cref{eq:efficient:tpr} at $\hat\alpha$ \\
\midrule
0.10 & 1 & 0.916 & 0.916 & 1.88 & 1.92 \\
0.10 & 2 & 0.856 & 0.894 & 2.67 & 2.69 \\
0.10 & 4 & 0.842 & 0.911 & 4.21 & 4.19 \\
0.10 & 8 & 0.409 & 0.803 & 4.10 & 4.37 \\
0.30 & 1 & 0.710 & 0.710 & 1.68 & 1.71 \\
0.30 & 2 & 0.516 & 0.623 & 2.00 & 2.01 \\
0.30 & 4 & 0.378 & 0.644 & 2.46 & 2.50 \\
0.30 & 8 & 0.213 & 0.633 & 2.62 & 2.68 \\
1.00 & 1 & 0.233 & 0.233 & 1.22 & 1.23 \\
1.00 & 2 & 0.116 & 0.195 & 1.22 & 1.23 \\
1.00 & 4 & 0.056 & 0.182 & 1.21 & 1.22 \\
1.00 & 8 & 0.037 & 0.226 & 1.28 & 1.29 \\
\bottomrule
\end{tabular}
\end{center}

\begin{figure}[t]
\centering
\begin{tikzpicture}
\begin{axis}[width=0.62\linewidth,height=5cm,xlabel={$\gamma$},ylabel={acceptance estimate},
  xmode=log,log basis x=2,xtick={1,2,4,8},xticklabels={1,2,4,8},ymin=0,ymax=1,
  grid=major,grid style={gray!20},legend pos=outer north east,legend style={font=\scriptsize}]
\addplot[sysblue,thick,mark=*] coordinates {(1,0.916) (2,0.894) (4,0.911) (8,0.803)};
\addlegendentry{$\hat\alpha$, noise 0.1}
\addplot[sysblue,thick,dashed,mark=o] coordinates {(1,0.916) (2,0.856) (4,0.842) (8,0.409)};
\addlegendentry{kept/proposed, noise 0.1}
\addplot[warnred,thick,mark=square*] coordinates {(1,0.710) (2,0.623) (4,0.644) (8,0.633)};
\addlegendentry{$\hat\alpha$, noise 0.3}
\addplot[warnred,thick,dashed,mark=square] coordinates {(1,0.710) (2,0.516) (4,0.378) (8,0.213)};
\addlegendentry{kept/proposed, noise 0.3}
\end{axis}
\end{tikzpicture}
\caption{Measured on CPU in this environment (toy models, 160 generated tokens per point).
The naive kept/proposed'' ratio falls with $\gamma$ although the per-position
acceptance $\hat\alpha$ is roughly constant: the bias described in the text.}
\label{fig:efficient:alpha-bias}
\end{figure}


Two things stand out. First, ``kept/proposed'' falls with $\gamma$ even when the
per-position acceptance does not, which is exactly the bias described above. Second,
the measured tokens per call track \cref{eq:efficient:tpr} evaluated at $\hat\alpha$,
within the noise of a 160-token run. Third, a poor draft (noise 1.0, $\hat\alpha\approx0.2$)
never gets past about 1.25 tokens per call however long it drafts, so every extra draft
step is pure overhead.

\subsection{Variants: where do drafts come from?}
\begin{itemize}
  \item \textbf{Separate small model} from the same family (as above). Simple; costs extra
        memory and its own KV cache; needs the same tokenizer.
  \item \textbf{Self-speculation / early exit} (e.g.\ LayerSkip, Elhoushi et al., 2024): run only
        the first $k$ layers of the target, then its head, as the draft. There are no
        extra weights, and the draft's KV for those layers can be reused. Our lab's
        \code{layer\_skip\_draft} builds this.
  \item \textbf{Extra decoding heads} (Medusa, Cai et al., 2024): $k$ small heads on top of
        the last hidden state predict tokens $t+2,\dots,t+k+1$ in parallel. Candidates are
        combined into a \emph{tree} and verified at once.
  \item \textbf{Feature-level drafting} (EAGLE, Li et al., 2024): a one-layer model
        autoregresses on the target's hidden states rather than on tokens, which yields
        much higher acceptance than an independent small LM. EAGLE-style drafts are the
        default in several production engines.
  \item \textbf{Prompt lookup / $n$-gram drafting}: propose the continuation of the most recent
        earlier occurrence of the last $n$ tokens. It is free, and remarkably effective for
        editing, retrieval-augmented answers and code, where outputs copy the input.
        Lookahead decoding (Fu et al., 2024) generates $n$-grams on the fly with Jacobi
        iterations.
\end{itemize}

\begin{figure}[t]
\centering
\begin{tikzpicture}[font=\scriptsize,
  n/.style={draw,rounded corners=1pt,minimum width=0.9cm,minimum height=0.45cm,fill=keygreen!20}]
\node[n,fill=sysblue!25] (root) at (0,0) {pending};
\node[n] (a) at (1.7,0.7) {cat};  \node[n] (b) at (1.7,-0.9) {dog};
\node[n] (c) at (3.4,1.2) {sat};  \node[n] (d) at (3.4,0.2) {is};
\node[n] (e) at (3.4,-0.9) {ran}; \node[n] (f) at (5.1,1.2) {on};
\foreach \u/\v in {root/a,root/b,a/c,a/d,b/e,c/f}{\draw[-Latex] (\u) -- (\v);}
% mask
\begin{scope}[xshift=6.8cm,yshift=1.5cm]
\def\names{{"p","cat","dog","sat","is","ran","on"}}
\foreach \r [count=\ri from 0] in {p,cat,dog,sat,is,ran,on}{
  \node[anchor=east] at (-0.1,-\ri*0.38-0.19) {\r};
  \node[rotate=60,anchor=west] at (\ri*0.38+0.19,0.1) {\r};}
% ancestors (row allowed cols): p:{p}; cat:{p,cat}; dog:{p,dog}; sat:{p,cat,sat}; is:{p,cat,is}; ran:{p,dog,ran}; on:{p,cat,sat,on}
\foreach \r/\c in {0/0,1/0,1/1,2/0,2/2,3/0,3/1,3/3,4/0,4/1,4/4,5/0,5/2,5/5,6/0,6/1,6/3,6/6}{
  \fill[sysblue!55] (\c*0.38,-\r*0.38) rectangle ++(0.38,-0.38);}
\draw[gray,step=0.38] (0,0) grid (2.66,-2.66);
\node[below,align=center] at (1.33,-2.75) {tree attention mask\\(row attends to its ancestors)};
\end{scope}
\end{tikzpicture}
\caption{Tree verification (Medusa, EAGLE, SpecInfer). Several candidate continuations
share prefixes. They are flattened into one block and verified in a single forward pass
with a mask that lets each node attend only to its ancestors and the committed context.
The longest accepted root-to-leaf path is kept.}
\label{fig:efficient:tree}
\end{figure}

\begin{systemslens}[Speculation and batching compete for the same slack]
Speculation converts idle compute into fewer sequential steps. A continuous-batching
server at high load has \emph{no} idle compute: its decode steps approach the ridge
point and verification of $\gamma+1$ tokens per sequence now costs real FLOPs. So
speculative decoding gives large latency wins at low load and can \emph{reduce}
throughput at high load. Production engines therefore adapt $\gamma$ to load or disable
speculation above a batch-size threshold. Decide with the roofline: if
$B(\gamma+1)$ tokens per step pushes arithmetic intensity past the ridge point, you are
paying for rejected drafts with real compute.
\end{systemslens}

% =====================================================================
\section{Splitting a model across GPUs for inference}\label{sec:efficient:parallel}

\Cref{ch:distributed} derived tensor, pipeline and expert parallelism for training.
Inference reuses the same partitionings, with different trade-offs because there is no
backward pass and decode moves tiny activations.

\paragraph{Tensor parallelism (TP).} Shard each attention block by heads and each MLP by
columns then rows (Megatron-style). Each layer needs two all-reduces of a
$B\times\dmodel$ activation (after attention and after the MLP). Every GPU reads only
$1/P$ of the weights and of the KV cache (KV heads are sharded too), so the bandwidth
floor drops by $P$. The price is $2L$ all-reduces per decode step, which are
latency-bound rather than bandwidth-bound at small $B$.

\begin{napkinmath}[Qwen3-32B with TP=2 on NVLink]
The weights are 65.5\,GB in bf16, so they do not fit one 80\,GB GPU with a useful KV
cache. With TP=2 each GPU holds 32.8\,GB, and the step floor falls from 19.6\,ms to
about 9.8\,ms. Per decode step and per sequence, the all-reduces move
$2\cdot 64 \cdot 5120 \cdot 2\,\text{B} \approx 1.3$\,MB, negligible in bandwidth. But
128 collectives at roughly 5--10\,$\mu$s each add 0.6--1.3\,ms of pure latency per
token. This is why TP is kept within an NVLink domain, and why engines fuse the
all-reduce with the following norm or use one-shot all-reduce kernels for small
messages.
\end{napkinmath}

\paragraph{Pipeline parallelism (PP).} Assign contiguous layers to GPUs. For a single
sequence this adds latency (the stages run one after another) but no speedup. It pays
off for throughput when many micro-batches are in flight, and it is the usual way to
span nodes connected by slower networks.

\paragraph{Expert parallelism (EP) for MoE.} A mixture-of-experts layer
(\cref{ch:modern}) activates only a few experts per token, so decode reads a fraction of
the weights. But the experts are spread across GPUs, and every MoE layer becomes two
all-to-all exchanges (dispatch tokens to experts, combine results). Large-EP serving
designs, such as DeepSeek's published inference setup, place one or a few experts per
GPU across many nodes and overlap the all-to-all with computation by splitting a batch
into two micro-batches. The new failure mode is \emph{load imbalance}: one hot expert
throttles the whole step. Systems therefore replicate popular experts and rebalance
placement from routing statistics.

\begin{figure}[t]
\centering
\begin{tikzpicture}[font=\scriptsize,
  gpu/.style={draw=sysblue,fill=sysbluebg,rounded corners=2pt,minimum width=2.3cm,minimum height=2.6cm},
  blk/.style={draw,fill=#1,minimum width=1.8cm,minimum height=0.42cm},
  ar/.style={draw=labpurple,fill=labpurple!20,rounded corners=1pt,minimum width=5.4cm,minimum height=0.36cm}]
\node[gpu] (g0) at (0,0) {}; \node[gpu] (g1) at (3,0) {};
\node[above] at (g0.north) {GPU 0}; \node[above] at (g1.north) {GPU 1};
\foreach \g/\x in {0/0,1/3}{
  \node[blk=fcompute!30] at (\x,0.85) {heads $\g$-half};
  \node[blk=fmemory!35] at (\x,0.3) {KV of those heads};
  \node[blk=fdata!30] at (\x,-0.45) {MLP cols/rows};
  \node[blk=fmemory!20] at (\x,-1.0) {$\tfrac12$ weights};}
\node[ar] at (1.5,1.55) {all-reduce after attention};
\node[ar] at (1.5,-1.65) {all-reduce after MLP};
% EP
\begin{scope}[xshift=7.2cm]
\foreach \i in {0,1,2,3}{\node[draw=keygreen,fill=keygreen!15,rounded corners=2pt,minimum width=1.0cm,minimum height=0.8cm] (e\i) at (\i*1.25,-0.6) {E\i};
  \node[draw,fill=sysblue!15,minimum width=1.0cm,minimum height=0.4cm] (t\i) at (\i*1.25,1.0) {tok \i};}
\foreach \a/\b in {0/2,1/0,2/2,3/1,0/3,2/1}{\draw[-Latex,fcomm,thick] (t\a.south) -- (e\b.north);}
\node[above=0.15cm of t1,xshift=0.6cm] {all-to-all dispatch (top-2 routing)};
\node[below=0.15cm of e2,xshift=-0.6cm,text=warnred] {E2 is hot $\Rightarrow$ straggler};
\end{scope}
\end{tikzpicture}
\caption{Left: tensor-parallel decode shards heads, KV cache and MLP across GPUs, with
two all-reduces per layer. Right: expert-parallel MoE layers route tokens to experts on
different GPUs with all-to-all exchanges; imbalance turns one expert into a straggler.}
\label{fig:efficient:parallel}
\end{figure}

% =====================================================================
\section{Structured output and constrained decoding}\label{sec:efficient:structured}

Applications often need output that \emph{parses}: JSON matching a schema, a function
call, a number in a \texttt{\textbackslash boxed\{\}} (\cref{ch:eval}). Rather than hoping,
we can \emph{constrain} sampling. Compile the grammar (regular expression, JSON schema
or context-free grammar) into an automaton over \emph{tokens}. At each step, set the
logits of tokens that would leave the language to $-\infty$, then sample as usual. The
model's probabilities are renormalised over the legal tokens.

\begin{lstlisting}[caption={Masking logits to a legal token set. The hard part, computing
\ident{allowed_ids} efficiently from a grammar state, is what libraries such as Outlines
and XGrammar provide.},label={lst:efficient:constrain}]
def constrained_step(logits, allowed_ids, temperature=0.0):
    mask = torch.full_like(logits, float("-inf"))
    mask[..., allowed_ids] = 0.0
    return sample_next_token(logits + mask, temperature)
\end{lstlisting}

The systems problem is cost. A BPE vocabulary of about 150k tokens must be checked
against the grammar at \emph{every} step, on the CPU, while the GPU waits. Outlines
(Willard \& Louf, 2023) precomputes, for each state of a regular expression's automaton,
the set of tokens that keep the automaton alive, which turns the per-step cost into an
index lookup. Context-free grammars (JSON with nesting) need a pushdown automaton, whose
stack makes the state space unbounded. XGrammar (Dong et al., 2024) splits tokens into
context-independent ones (checkable ahead of time) and a small context-dependent
remainder checked at runtime, and overlaps mask computation with the GPU forward pass.
Constraints also interact with tokenization: forcing \texttt{"\}"} as a single token
when the model would naturally emit \texttt{"\}\textbackslash n"} pushes it
off-distribution. This is the ``token healing'' problem, and it is why good libraries
allow every token whose \emph{characters} are consistent with the grammar.

\begin{keyidea}
Constrained decoding changes the output distribution: it samples from $p$ restricted to
the language, not from $p$ conditioned on producing a valid string. For short structured
fields the difference is negligible. For long reasoning followed by an answer, let the
model think unconstrained and constrain only the final answer span.
\end{keyidea}

% =====================================================================
\section{The serving stack around the engine}\label{sec:efficient:stack}

An inference engine (\cref{ch:serving}) turns token ids into token ids. A service must
also accept HTTP requests, apply chat templates, tokenize, schedule, stream text back,
and survive load. \Cref{fig:efficient:lifecycle} traces one request through an
OpenAI-compatible server such as those shipped with vLLM, SGLang, TensorRT-LLM or TGI.

\begin{figure}[t]
\centering
\begin{tikzpicture}[font=\scriptsize,
  box/.style={draw,rounded corners=2pt,minimum height=0.9cm,align=center,fill=#1},
  arr/.style={-Latex,thick}]
\node[box=gray!10] (client) at (0,0) {client\\(SDK, curl)};
\node[box=sysblue!12] (api) at (2.4,0) {HTTP/API\\server};
\node[box=sysblue!12] (tmpl) at (4.7,0) {chat template\\+ tokenizer};
\node[box=fcompute!25] (sched) at (7.1,0) {scheduler\\(cont.\ batching)};
\node[box=fcompute!25] (eng) at (9.6,0) {GPU engine\\(paged KV)};
\node[box=sysblue!12] (detok) at (7.1,-1.8) {incremental\\detokenizer};
\node[box=sysblue!12] (sse) at (2.4,-1.8) {SSE stream\\\texttt{data: \{...\}}};
\draw[arr] (client) -- (api); \draw[arr] (api) -- (tmpl); \draw[arr] (tmpl) -- (sched);
\draw[arr] (sched) -- (eng);
\draw[arr] (eng.south) |- (detok.east) node[pos=0.25,right] {token ids / step};
\draw[arr] (detok) -- (sse) node[midway,above] {text deltas};
\draw[arr] (sse.west) -| (client.south);
\draw[decorate,decoration={brace,amplitude=4pt,mirror}] (3.6,-2.5) -- (10.6,-2.5)
  node[midway,below=4pt,align=center] {CPU work that must stay off the GPU's critical path\\
  (separate processes in production engines)};
\node[above=0.1cm of sched,text=srcgray] {admission, priorities, preemption};
\node[above=0.1cm of api,text=srcgray] {auth, rate limits};
\end{tikzpicture}
\caption{Life of a streaming chat request. Tokenization, detokenization and HTTP handling
run on the CPU. Production engines move them into separate processes so a slow client or
a long prompt cannot stall the batched GPU loop.}
\label{fig:efficient:lifecycle}
\end{figure}

\paragraph{Streaming.} Responses are streamed as server-sent events (SSE), one small JSON
delta per step. This decouples \emph{time to first token} (TTFT) from total latency.
Two details matter. First, detokenization must be \emph{incremental and UTF-8 aware}:
decoding one BPE token in isolation can split a multi-byte character (an emoji is 4
bytes and may span several tokens). The usual fix decodes a small window of recent ids
and emits only text that no longer ends in the replacement character U+FFFD. Our chat app
uses exactly this (\code{IncrementalDetokenizer}). Second, the CPU threads that
serialise and send deltas must never block the GPU loop. In Python servers that means a
separate process or thread for the engine, communicating through queues.

\paragraph{Load testing.} To measure a serving system, drive it with an \emph{open-loop}
workload: requests arrive by a Poisson process at rate $\lambda$, independent of how fast
the server responds. Closed-loop clients (``$N$ users each waiting for their reply'')
self-throttle and hide queueing collapse. Sweep $\lambda$. At each rate record the
percentiles (p50/p90/p99) of TTFT and inter-token latency (ITL), and the \emph{goodput}:
completed requests per second that met their latency SLOs. Use realistic length
distributions (prompt and output lengths sampled from logs or from datasets such as
ShareGPT), because the prefill/decode mix determines the result.

\begin{napkinmath}[Little's law sizes your KV cache]
Little's law says $\bar N = \lambda\,\bar W$: the mean number of requests in the system
equals arrival rate times mean time in the system. At $\lambda=20$ requests/s with
replies that take $\bar W=12$\,s, about 240 sequences are resident on average. At 3000
tokens each and 147.5\,KB per token (Qwen3-8B), that is $240\cdot 3000\cdot 147.5$\,KB
$\approx 106$\,GB of KV cache, more than an 80\,GB GPU. The server must shard, quantize
the KV, or queue, and queueing raises $\bar W$, which raises $\bar N$ further. This
feedback loop is why latency explodes near capacity.
\end{napkinmath}

\begin{systemslens}[Disaggregation, revisited]
Prefill is compute-bound and bursty; decode is memory-bound and steady. Running them on
the same GPUs makes long prompts inflate the ITL of everyone else, and chunked prefill
(\cref{ch:serving}) only mitigates this. Disaggregated serving sends prefill to one pool,
ships the KV cache over NVLink, InfiniBand or RDMA, and decodes in another pool, so each
pool can use its own parallelism and even its own hardware generation. The KV transfer
is the new bottleneck. For Qwen3-8B a 4k-token prompt produces about 0.6\,GB of KV,
roughly 12\,ms at 50\,GB/s.
\end{systemslens}

% =====================================================================
\section{A chat interface for our models}\label{sec:efficient:chat}\covers{RM G}

We now put a browser front end on the from-scratch Qwen3 implementation. It is useful
for poking at the checkpoints we train later (instruction-tuned in \cref{ch:instruct},
RL-trained in \cref{ch:grpo}, distilled in \cref{ch:distill}), and it is a compact case
study of everything above: chat templates, streaming, detokenization, event loops and KV
reuse. We use Chainlit, an open-source Python package that provides a ChatGPT-like web
UI and calls our Python functions when messages arrive.

\subsection{Installing Chainlit}\covers{RM G.1}
Chainlit is a regular PyPI package:
\begin{lstlisting}[style=shell]
pip install chainlit          # or, with uv:  uv add chainlit
\end{lstlisting}
It is deliberately \emph{not} in the lab's \code{requirements.txt}: nothing else in the
book depends on it, and it pulls in a web stack (an ASGI server, websockets).

\subsection{Running an app as a script}\covers{RM G.2}
A Chainlit app is an ordinary Python file that Chainlit \emph{imports} and inspects
for decorated handlers: \code{@cl.on\_chat\_start} runs when a browser session opens,
and \code{@cl.on\_message} runs for every user message. You start the server with
\begin{lstlisting}[style=shell]
chainlit run chat_single_turn.py      # serves http://localhost:8000
\end{lstlisting}
Because \code{chainlit run} owns the process's command line, our scripts cannot use
\code{argparse}. Configuration therefore comes from environment variables
(\code{QWEN3\_DIR}, \code{CHECKPOINT\_PATH}, \code{MAX\_NEW\_TOKENS},
\code{THINKING}, \code{QUANT\_BITS}), read once at import time in \code{chat\_common.py}. Module import
also happens exactly once per server process, so that is where the model is loaded and
shared by all browser sessions.

\subsection{Where the scripts live}\covers{RM G.3}
Everything is in \ident{code/labs/lab20_quant_specdec/}: \code{chat\_common.py} (model
loading, streaming, detokenization, history trimming), \code{chat\_single\_turn.py} and
\code{chat\_multi\_turn.py}. They import only \code{llmsys} and the lab's
\code{solution.py} (for optional quantized serving and prefix matching), so there is
nothing to download beyond the Qwen3 weights.

\subsection{The single-turn app}\covers{RM G.4}
\Cref{lst:efficient:single} is the whole single-turn app. Each message is wrapped in the
ChatML template, tokenized, and streamed back token by token into an initially empty
message bubble.

\begin{lstlisting}[caption={\ident{chat_single_turn.py}: every message is answered
independently.},label={lst:efficient:single}]
import chainlit as cl
import torch
from chat_common import (DEVICE, THINKING, IncrementalDetokenizer,
                         load_model_and_tokenizer, stop_token_ids, stream_tokens)

MODEL, TOKENIZER = load_model_and_tokenizer()   # once per server process
STOP = stop_token_ids(TOKENIZER)                # <|im_end|>, <|endoftext|>

@cl.on_chat_start
async def on_start():
    await cl.Message(content="Single-turn mode.").send()

@cl.on_message
async def on_message(message: cl.Message):
    prompt = TOKENIZER.apply_chat_template(message.content, enable_thinking=THINKING)
    ids = torch.tensor([TOKENIZER.encode(prompt)], device=DEVICE)
    reply = cl.Message(content="")
    await reply.send()                           # empty bubble to stream into
    detok = IncrementalDetokenizer(TOKENIZER)
    async for tok in stream_tokens(MODEL, ids, STOP):
        if tok in STOP:
            break
        piece = detok.push(tok)
        if piece:
            await reply.stream_token(piece)
    await reply.update()
\end{lstlisting}

Two systems details hide in \code{chat\_common.py}. First, \code{load\_model\_and\_tokenizer}
loads the official snapshot by default, or, if \code{CHECKPOINT\_PATH} is set, builds a
fresh \code{Qwen3Model} and loads your own state dict. The tokenizer always comes from
the snapshot directory, so it must match the checkpoint: base and reasoning variants
use different chat conventions, and a checkpoint fine-tuned on one should be served
with the same template. Second, \code{stream\_tokens} runs generation in a worker thread
and passes tokens to the asyncio event loop through a queue
(\cref{lst:efficient:stream}). Calling the model directly inside the \code{async} handler
would block the event loop for the whole reply: the UI would not repaint and every
other session would freeze.

\begin{lstlisting}[caption={Bridging blocking generation and an async web framework.},
label={lst:efficient:stream}]
async def stream_tokens(model, input_ids, stop_ids, cache=None):
    loop = asyncio.get_running_loop()
    queue, done = asyncio.Queue(), object()
    def worker():                                 # runs in a plain thread
        try:
            for tok in generate_stream(model, input_ids, MAX_NEW_TOKENS, cache=cache):
                t = int(tok[0, 0])
                loop.call_soon_threadsafe(queue.put_nowait, t)
                if t in stop_ids:
                    break
        finally:
            loop.call_soon_threadsafe(queue.put_nowait, done)
    threading.Thread(target=worker, daemon=True).start()
    while (item := await queue.get()) is not done:
        yield item
\end{lstlisting}

\subsection{Running it, with your own checkpoints}\covers{RM G.5}
\begin{lstlisting}[style=shell]
huggingface-cli download Qwen/Qwen3-0.6B --local-dir qwen3-0.6b
QWEN3_DIR=qwen3-0.6b chainlit run chat_single_turn.py
# serve a checkpoint from Chapter 27 or 20 instead, 4-bit quantized:
QWEN3_DIR=qwen3-0.6b CHECKPOINT_PATH=runs/grpo/step500.pth QUANT_BITS=4 \
    chainlit run chat_single_turn.py
\end{lstlisting}
Chainlit usually opens a browser tab. Otherwise visit \url{http://localhost:8000}. With
\code{THINKING=1} the reasoning model first emits a \texttt{<think>\,\dots\,</think>}
block; a simple prompt such as ``What is 1+1?'' can trigger several hundred tokens of
deliberation, which makes \code{MAX\_NEW\_TOKENS} a real cost knob
(\cref{ch:scaling}).

\subsection{Multi-turn conversations}\covers{RM G.6}
\paragraph{What multi-turn means.} The single-turn app \emph{displays} a conversation but
the model sees only the latest message. Ask ``What is 1+1?'' and then ``What did I just
ask?'' and the model has no idea. A multi-turn chat re-sends the history with every
request, because the model itself is stateless: the ``memory'' of a chat assistant is
nothing but its prompt.

\paragraph{The multi-turn app.} \code{chat\_multi\_turn.py} keeps a per-session list of
messages in Chainlit's per-session store (\code{user\_session}), seeded with a system
message. On each turn it
appends the user message, serialises the whole history with the ChatML template
(\cref{fig:efficient:chatml}), generates, and appends the assistant reply.

\begin{figure}[t]
\centering
\resizebox{\linewidth}{!}{%
\begin{tikzpicture}[font=\scriptsize\ttfamily,
  seg/.style={draw,rounded corners=1pt,align=left,anchor=west,inner sep=3pt,fill=#1,text width=9.4cm},
  lab/.style={font=\scriptsize\sffamily,text=srcgray,anchor=east}]
\node[seg=gray!12] (s) at (0,0) {<|im\_start|>system\textbackslash n You are a helpful assistant.<|im\_end|>};
\node[seg=sysblue!12] (u1) at (0,-0.55) {<|im\_start|>user\textbackslash n What is 1+1?<|im\_end|>};
\node[seg=keygreen!15] (a1) at (0,-1.1) {<|im\_start|>assistant\textbackslash n 2<|im\_end|>};
\node[seg=sysblue!12] (u2) at (0,-1.65) {<|im\_start|>user\textbackslash n What did I just ask?<|im\_end|>};
\node[seg=keygreen!35] (a2) at (0,-2.2) {<|im\_start|>assistant\textbackslash n \textit{(model continues here)}};
\node[lab] at (-0.1,0) {system}; \node[lab] at (-0.1,-0.55) {turn 1};
\node[lab] at (-0.1,-1.1) {reply 1}; \node[lab] at (-0.1,-1.65) {turn 2};
\node[lab] at (-0.1,-2.2) {header};
\draw[decorate,decoration={brace,amplitude=5pt},thick,sysblue]
  (10.05,0.25) -- (10.05,-1.35) node[midway,right=6pt,font=\scriptsize\sffamily,align=left]
  {identical to last turn's tokens:\\K/V already cached $\Rightarrow$ skip};
\draw[decorate,decoration={brace,amplitude=5pt},thick,warnred]
  (10.05,-1.4) -- (10.05,-2.45) node[midway,right=6pt,font=\scriptsize\sffamily,align=left]
  {new suffix: prefill only this};
\end{tikzpicture}}
\caption{The multi-turn prompt is the whole history serialised in ChatML. Because each
turn's prompt extends the previous one, the session can keep its KV cache and prefill
only the new suffix.}
\label{fig:efficient:chatml}
\end{figure}

\paragraph{Fitting the context window.} Histories grow without bound, so before each
turn we must leave room for \code{MAX\_NEW\_TOKENS}. The simplest strategy left-truncates
the token sequence. That can cut a message in half and even drop the
\texttt{<|im\_start|>} marker, producing a prompt shape the model never saw in
training. Our \code{trim\_history} instead drops the oldest \emph{whole} user/assistant
turns and never drops the system message. More elaborate strategies summarise old turns
with the model itself.

\paragraph{How history meets the KV cache.} Re-sending the history means re-prefilling
it, an $O(T)$ cost that grows every turn, so a 20-turn chat does quadratic total work.
The app therefore keeps the session's \code{KVCache} together with the token ids it
covers. Each turn it computes the longest common prefix between the cached ids and the
new prompt (\code{common\_prefix\_len}), truncates the cache to that length, and
prefills only the suffix (\cref{lst:efficient:multiturn}). This is a single-session
version of the prefix caching that engines such as vLLM and SGLang perform across all
requests (\cref{ch:serving}).

\begin{lstlisting}[caption={The heart of \ident{chat_multi_turn.py}: history in, KV cache
reused.},label={lst:efficient:multiturn}]
history.append({"role": "user", "content": message.content})
history = trim_history(history, TOKENIZER, CONTEXT - MAX_NEW_TOKENS)
prompt_ids = TOKENIZER.encode(TOKENIZER.apply_chat_template(history, THINKING))
cache, cached = cl.user_session.get("cache"), cl.user_session.get("cached_ids")
reuse = min(common_prefix_len(cached, prompt_ids), len(prompt_ids) - 1)
truncate_cache(cache, reuse)                     # drop anything that diverged
new_ids = torch.tensor([prompt_ids[reuse:]], device=DEVICE)
...                                              # stream as in single-turn
cl.user_session.set("cached_ids", (prompt_ids + generated)[: cache.seq_len(0)])
history.append({"role": "assistant", "content": reply.content})
\end{lstlisting}

\begin{pitfall}[When the prefix silently stops matching]
Prefix reuse is only valid if re-serialising the history reproduces \emph{exactly} the
tokens that were cached. Three things break it. (1) Templates that rewrite history:
Qwen3's official template removes earlier turns' \texttt{<think>} blocks, so the
re-serialised assistant turn is shorter than what was generated. Our simplified
template keeps them, so the prefix holds, but the model sees its old reasoning (more
tokens, and different from its training distribution). (2) Tokenization is not
compositional: the text of a reply, re-encoded in context, can tokenize differently
from the ids that were sampled. (3) Trimming old turns changes the prefix from the
start. \code{common\_prefix\_len} makes all three cases \emph{safe}, since it falls back
to recomputation, but they silently cost performance. Log the reuse ratio (the app
prints it after every reply).
\end{pitfall}

\paragraph{Recommendations.} Use the multi-turn app with models trained on multi-turn chat
data, such as the official Qwen3 reasoning checkpoint. Checkpoints trained only on
single prompt/response pairs, such as the distilled models of \cref{ch:distill} or our
small instruction-tuned models, were never taught to condition on earlier assistant
turns. Inspect those with the single-turn app; the multi-turn app tends to confuse them.

% =====================================================================
\section{Summary}
\begin{itemize}
  \item A decode step costs, to first order, the bytes of weights plus KV cache divided
        by HBM bandwidth (\cref{eq:efficient:bytes}). Decide which term dominates before
        choosing an optimisation.
  \item Uniform quantization trades resolution for bytes; outliers are the enemy, and
        finer granularity (per-channel, per-group, per-block microscaling) is the main
        weapon. Weight-only int8/int4 speeds up memory-bound decode; W8A8/FP8 speeds up
        compute-bound prefill and large batches; KV quantization raises both speed and
        capacity in long-context serving.
  \item GPTQ uses second-order information to compensate rounding error; AWQ and SmoothQuant
        rescale channels so that what matters lands on the grid.
  \item Speculative decoding is \emph{lossless}: accept with $\min(1,p/q)$, resample from
        $\max(p-q,0)$. Expected tokens per round are $(1-\alpha^{\gamma+1})/(1-\alpha)$;
        its benefit shrinks as batching consumes the idle compute it relies on.
  \item Tensor parallelism divides the bandwidth floor at the cost of $2L$ latency-bound
        collectives per step; expert parallelism trades all-to-alls and load balance for
        sparse compute.
  \item A service adds templates, streaming, incremental detokenization, async I/O,
        open-loop load testing and SLO-based goodput around the engine.
  \item A chat UI is a thin layer: the model is stateless, memory is the prompt, and the
        KV cache can make re-sending history nearly free.
\end{itemize}

% =====================================================================
\section{Exercises}
\begin{exercise}[bytes per token]
Using \cref{eq:efficient:bytes}, find the batch size at which Qwen3-0.6B's KV cache bytes
equal its weight bytes at 2048 tokens of context (bf16 for both). What does this say
about weight-only quantization for small models in a busy server?
\end{exercise}
\begin{exercise}[rounding noise]
Show that if $x$ is uniformly distributed over a quantization bin of width $s$, the mean
squared rounding error is $s^2/12$. Use this to derive the ``6\,dB per bit'' rule for a
uniform quantizer of a signal spanning its full range.
\end{exercise}
\begin{exercise}[outliers]
A weight row has 127 values drawn from $\mathcal N(0,1)$ and one value equal to 40.
Estimate the fraction of the 127 normal values that round to zero under per-row
symmetric int4 quantization, and under group-wise int4 with $g=32$.
\end{exercise}
\begin{exercise}[metadata]
Compute the effective bits per weight of: (a) int4 with $g=64$, fp16 scale and fp16
zero; (b) MXFP4 (blocks of 32, one 8-bit shared scale); (c) NVFP4 (blocks of 16, one
FP8 scale, ignoring the per-tensor scale).
\end{exercise}
\begin{exercise}[GPTQ derivation]
Starting from \cref{eq:efficient:gptq} for a single row $\vect w$, show that if column
$i$ is fixed to $q_i$, the optimal update of the remaining columns is
$\delta\vect w_{-i} = -\frac{w_i-q_i}{[H^{-1}]_{ii}}[H^{-1}]_{-i,i}$, and that the
increase in loss is $(w_i-q_i)^2/[H^{-1}]_{ii}$.
\end{exercise}
\begin{exercise}[SmoothQuant]
Show that \cref{eq:efficient:smooth} leaves $XW^\top$ unchanged, and that with
$\alpha=0.5$ the transformed per-channel maxima satisfy
$\max|X'_{:,j}|=\max|W'_{:,j}|$. How would you fold $\mathrm{diag}(\vect s)^{-1}$ into
a preceding RMSNorm?
\end{exercise}
\begin{exercise}[lossless speculation]
Extend the proof of \cref{prop:efficient:exact} to two consecutive positions: show that
the pair of tokens emitted when the first draft is accepted and the second is
verified has distribution $p(x_1)\,p(x_2\mid x_1)$.
\end{exercise}
\begin{exercise}[optimal draft length]
For $c=0.05$, find numerically the $\gamma$ that maximises \cref{eq:efficient:speedup}
for $\alpha\in\{0.6,0.8,0.9\}$. Derive the condition for $\gamma\to\gamma+1$ to help.
\end{exercise}
\begin{exercise}[speculation under load]
A server runs Qwen3-8B on one H100 at batch 128 with 2k-token contexts. Using the
roofline, estimate whether verifying $\gamma+1=5$ tokens per sequence keeps the step
memory-bound. What $\gamma$ would you choose?
\end{exercise}
\begin{exercise}[TP latency]
Model a tensor-parallel decode step as $t = t_{\text{mem}}/P + 2L\,(a + m/\beta)$ with
per-collective latency $a$, message size $m$ and link bandwidth $\beta$. For Qwen3-32B
with $a=8\,\mu$s, find the $P$ that minimises per-token latency at batch 1.
\end{exercise}
\begin{exercise}[load testing]
Explain why a closed-loop benchmark with 64 concurrent users can report excellent p99
latency for a server that collapses under an open-loop load of the same average rate.
\end{exercise}
\begin{exercise}[chat memory]
In the multi-turn app, after 20 turns of 150 tokens each, how many prompt tokens are
prefilled in total with and without KV reuse? Which of the three pitfalls in
\cref{sec:efficient:chat} would you expect to hit first with the official Qwen3
template, and how would you fix it?
\end{exercise}

% =====================================================================
\section{Lab 20: Quantize, speculate, and serve}\label{sec:efficient:lab}
\begin{labbox}{20}{Quantize, speculate, and serve}
\labdir{lab20_quant_specdec}
\labsection{Goals}
Implement weight-only int8 and int4 quantization from the bit level up, build a
speculative decoder with KV-cache rollback that is provably lossless, and serve a model
through a streaming multi-turn chat UI with cross-turn KV reuse.

\labsection{Setup}
\code{cd code; LAB\_IMPL=starter python -m pytest -q labs/lab13\_quant\_specdec}. All
tests run on CPU with tiny random Qwen3/GPT models. \code{bench.py} contains the
real-model experiments (Qwen3-0.6B/1.7B snapshots), and the chat apps need
\code{pip install chainlit}.

\labsection{Tasks}
\begin{enumerate}
  \item \code{quantize\_absmax} / \code{dequantize\_absmax}: symmetric int8, per-tensor or
        per-channel. The test injects an outlier row; per-channel must win.
  \item \code{quantize\_groupwise} / \code{dequantize\_groupwise}: asymmetric 4-bit groups
        with zero points. Then \code{pack\_int4} / \code{unpack\_int4} (two codes per byte).
  \item \code{quantize\_linear\_layers}: deep-copy a model and replace every
        \code{nn.Linear} except the head with \code{WeightOnlyLinear}. Check relative
        output error ($<1\%$ for int8, $<10\%$ for int4) and the byte savings.
  \item \code{accept\_or\_resample}: the speculative sampling rule. The test checks by
        Monte Carlo that the output follows $p$.
  \item \code{verify\_drafts}: accept a prefix and emit one extra token (residual or bonus).
  \item \code{speculative\_generate}: the full loop with both KV caches, rollback, and
        the ``draft never saw $d_\gamma$'' fix. Greedy output must equal plain greedy
        decoding of the target, using \code{layer\_skip\_draft} as the draft.
  \item \code{expected\_tokens\_per\_round}, and \code{common\_prefix\_len} for the chat app.
\end{enumerate}

\labsection{Measurements}
\begin{enumerate}
  \item \code{python bench.py toy}: reproduce the acceptance table of
        \cref{sec:efficient:spec}. Plot measured tokens per target call against
        \cref{eq:efficient:tpr} evaluated at $\hat\alpha$, for each noise level.
  \item With real checkpoints: \code{bench.py quant} reports KL to the bf16 reference and
        linear-layer memory for int8 and int4. Expect int8 to be nearly lossless and int4 to
        show a measurable but small KL. Then time decode with your PyTorch
        \code{WeightOnlyLinear} and note that it is \emph{slower} than bf16. Explain why
        with \cref{lst:efficient:quantlinear} and fix it with the fused kernel of Lab~A.
  \item \code{bench.py spec} with Qwen3-0.6B drafting for Qwen3-1.7B (or 4B): tokens/s and
        $\hat\alpha$ for $\gamma\in\{2,4,6,8\}$ on a math prompt and a creative prompt, at
        temperature 0 and 0.7.
  \item Run \code{chat\_multi\_turn.py} for ten turns and record the reuse ratio it prints.
        Switch \code{THINKING} on and off and explain the difference.
\end{enumerate}

\labsection{Deliverables}
Passing tests; a one-page report with the acceptance plot, the quantization KL/memory
table, the speculative speedup table, and an answer to ``at what batch size would you
turn speculation off for this pair of models, and why?''

\labsection{Stretch goals}
\begin{itemize}
  \item Implement \code{gptq\_quantize} for a real layer of Qwen3-0.6B, using 128
        calibration sequences, and compare layer output error with RTN at int4 and int3.
  \item Implement tree verification: draft the top-2 tokens at each of 2 levels (a tree of
        6 nodes), build the ancestor mask of \cref{fig:efficient:tree}, and verify with one
        forward pass. (You will need to pass a custom mask into attention.)
  \item Add int8 KV-cache quantization (\code{llmsys.quant.quantize\_kv}) to a copy of
        \code{KVCache} and measure the KL it induces at 4k context.
  \item Write a tiny OpenAI-compatible \code{/v1/chat/completions} endpoint with SSE
        streaming (FastAPI), and load-test it open-loop at increasing Poisson rates.
\end{itemize}
\end{labbox}
